\section{Machine checks}\label{app:checks}

Every finite quantity used in the proofs is recomputed by programs in the archive \cite{BB26code}, and the combinatorial and arithmetic statements among them are also proved in Lean~4. No statement of the paper rests on these checks; they guard against slips in the hand computations.

\paragraph{Lean.} The file \path{lean/HodgeClosure.lean} (Lean~4.24.0, no Mathlib) contains thirty-six theorems. Section~1 encodes the eight statements and the rule set of Table~\ref{tab:rules} as data, computes closures by iteration, and proves by evaluation that the closures of $\{\mathrm{F2},\mathrm{F3}'\}$, $\{\mathrm{F3}\}$, $\{\mathrm{L},\mathrm{M}\}$, $\{\mathrm{L},\mathrm{F3}'\}$ and $\{\mathrm{V},\mathrm{F3}'\}$ contain $\HC$, that the closures of $\{\mathrm{F2}\}$, $\{\mathrm{F3}'\}$, $\{\mathrm{L}\}$, $\{\mathrm{M}\}$, $\{\mathrm{V}\}$, $\{\AB\}$, $\{\mathrm{L},\mathrm{V}\}$ and $\{\mathrm{M},\mathrm{F3}'\}$ do not, and that the list of minimal sufficient sets, computed from the rule set alone, has exactly six members, those of Proposition~\ref{prop:minimal}. Section~2 works in $\ZZ[\sqrt{-d}]$ and proves that the hermitian form $\eta$ transforms under $\tau$ by $N(\tau)$, that the norm is multiplicative, that $\eta$ is not $K^{\times}$-invariant, and that the characters $N(\tau)^{n}$ and $\tau^{2n}$ differ at explicit $\tau$ (Section~\ref{sec:earlier}). Section~3 proves the dimension counts of Section~\ref{sec:f2}: $3n<n^{2}$ for $n\ge4$, $m^{2}+m'^{2}<(m+m')^{2}$ for $m,m'\ge1$ with $2mm'$ mixed directions, the values $r=57$ and $104$ of the lower bound, and the Hodge class count $26$ of Theorem~\ref{thm:link}(i) as a sum of K\"unneth contributions. Section~4 proves the polynomial identities behind the Fricke families of Section~\ref{ssec:fricke}: $c_{4}^{3}-c_{6}^{2}=1728\,t\,(t^{2}+13t+49)^{2}$ at level $7$ and $1728\,t\,(t^{2}+22t+125)^{3}$ at level $5$, by exact arithmetic on coefficient lists, together with the root products $7^{2}$ and $5^{3}$, the Euler numbers $12$ and $24$, the Shioda--Tate counts giving Picard number at least $16$, the transcendental rank $6$, and the dimensions $(n+1)/2$ of the dihedral Hecke algebras. Every theorem is checked by \texttt{\#print\,axioms}: thirty-four use no axiom, and the two inequalities proved with the arithmetic tactic use \texttt{propext} and \texttt{Quot.sound}.

The Lean certificate of \cite{BMB26}, \texttt{HodgeObstruction.lean} (Lean~4.34.0, \leancount{} theorems, each axiom-free or depending on \texttt{propext} only), carries the finite computations of that paper; it is cited and not repeated here. The reduction of Section~\ref{sec:reduction} uses only the hand-proved parts of \cite{BMB26}.

\paragraph{Python.} \path{code/sqrt6_chain.py} (sympy) implements Hilbert symbols and Hasse invariants exactly, verifies the product formula on random pairs, classifies $U_{\QQ}$, $E_{8}(-1)_{\QQ}$ and the invariant lattices, reproduces Varesco's embedding results for $\sqrt2$ and $\sqrt3$ as controls, searches the $165$ sign-admissible small forms $q_{E}$ over $\QQ(\sqrt6)$ ($98$ embed in $V_{3}$, $67$ do not), exhibits the complements $\langle1,-10\rangle$ and $\langle1,-1,-1,-1,-3,-30\rangle$, and checks both orders of the chain: twenty-five checks. \path{code/fricke_families.py} computes the fibre types, Euler numbers and Picard bounds of the Fricke curves and their twists at levels $5$ and $7$, checks numerically that the roots of $j(t)=j(\tau)$ are carried by $t\mapsto N^{k}/t$ to the $j$-invariants of the $N$-isogenous curves, verifies the twist identity $u_{p,q}(w(t))=pq\,u_{w(p),w(q)}(t)/t^{2}$, the M\"obius transformations preserving the labelled points, the embedding of every rank-six form of signature $(2,4)$ in $V_{2}$ and $V_{3}$ after any scaling, and the Hecke algebras of $(D_{n},\langle\tau\rangle)$ for $n=5,7,9,11$: forty-one checks. \path{code/product_object.py} builds the exterior-algebra model of $H^{*}(A\times A')$ for $(m,m')=(2,1),(2,2)$, the tangent space of $\cD_{n,n}$ and of the product locus, the first-order Hodge loci of the exterior products and of the object of Theorem~\ref{thm:product}, the Vandermonde determinant, the vanishing of the cross terms, and the Hankel rank and $r(\gamma)$ of the final class; \path{code/convolution_checks.py} computes in the same model the rank of $v\mapsto v\lrcorner\theta_{A}$ on $T\cD_{n,n}$, the contraction ranks of $\ch(\mathcal{E}')$ and $\ch(\mathcal{E}'')$ in Proposition~\ref{prop:convolution}(iii), and the trace of an upper triangular times a strictly upper triangular block matrix. In these two programs ranks are computed in floating point with a relative tolerance of $10^{-7}$ and agree with the closed formulas. \path{code/sixfold_link.py} checks the alternation and invariances of $\psi_{c}$, the decomposition of the $K^{1}$-invariants, the ratio of the induced forms at the three places, the symmetry of $T_{X}(-1)$, the count $26$ and the weight support of $\xi$; \path{code/link_weights.py} computes the Hodge classes of degree four on $X\times X\times Z_{K}$ under the three groups of Section~\ref{ssec:f2stop2} and the vanishing of $\Hom(X_{c},Z_{K,c})$ at the compositum CM point. \path{code/closure_graph.py} recomputes the closures and the minimal sets, and \path{code/bypass_checks.py} the Horn closure with the facts of the literature on abelian varieties adjoined and the blow-up bookkeeping of Proposition~\ref{prop:bypass}(ii).

\paragraph{Julia, C, C\texttt{++} and shell.} \path{code/closure_graph.jl}, \path{code/closure_graph.c} and \path{code/closure_graph.cpp} recompute the closures of all $256$ subsets of the eight statements and the six minimal sets by three independent implementations of the fixed-point iteration; \path{code/eigencharacter.jl} recomputes the characters of Section~\ref{sec:earlier} in exact integer arithmetic, and \path{code/hilbert.c} the Hilbert symbols and the invariants of Table~\ref{tab:invariants}; \path{code/run_all.sh} runs every program and the Lean file and prints a line per check. The expected outputs are recorded in \path{code/EXPECTED.md}.
